\documentclass{gatech-thesis}
% The default options:
%   \documentclass[11pt,letterpaper,oneside,%
%       doublespaced,normalmargins,final]{gatech-thesis}
% will generate a document that conforms to the graduate studies office
% guidelines.

 \ifx\pdfoutput\undefined
   \usepackage[dvips,final]{graphicx}  % using latex+dvips
   \usepackage[dvips,usenames]{color}
 \else
   \usepackage[pdftex,final]{graphicx} % using pdflatex
   \usepackage[pdftex,usenames]{color}
 \fi

%define stuff in preamble
 \degree{Doctor of Philosophy}
 \department{School of Electrical and Computer Engineering}
 \title{20,000 Leagues Under the Sea \protect\\ by Jules Verne}
 \author{Frodo Baggins}
 \principaladvisor{Elrond Halfelven}
 \committeechair{Gandalf Ol\'orin}
 \firstreader{Samwise Gamgee}
 \secondreader{Peregrin Took}[School of Industrial and Systems Engineering]
 \thirdreader{Meriadoc Brandybuck}[Department of Combinatorics and Optimization][University of Waterloo]
 \submitdate{August 2002}
 \approveddate{21 August 2002}
 \copyrightyear{2002} %add one if thesis submitted in Dec.

% \thesisproposalfalse       % default
% \titlepagetrue             % default
% \signaturepagetrue         % default
% \copyrightfalse            % default
% \figurespagetrue           % default
% \tablespagetrue            % default
% \contentspagetrue          % default
% \dedicationheadingfalse    % default
% \bibpagetrue               % default
% \strictmarginstrue         % default

\bibfiles{jules-verne-bib}

\begin{document}
\bibliographystyle{gatech-thesis}
\setchaptertocdepth{2}
%
\begin{preliminary}
\begin{dedication}
\begin{center}
  A dedication would go here...
\end{center}

\end{dedication}
\begin{acknowledgements}
\begin{acknowledgements}
Acknowledgement text goes here.  You could also (optionally)
have a forward or a preface, instead of an acknowledgement.

The text used as ``filler'' for Chapters 1, 2, 3, and 
Appendix A and B in this example was taken from the Project
Gutenberg edition of Jules Verne's \textit{20,000 Leagues Under the Sea.}
That text can be freely distributed according to the Gutenberg
license, in any markup language desired (here, \LaTeX) --- so long
as the original text version is also available.  To abide by
these requirements, the \textbf{full} original etext of
\textit{20,000 Leagues} is included in this archive as 
\programname{2000010.txt}.

Check \programname{2000010.txt} for Gutenberg licence and distribution
information.

This class file, \programname{gatech-thesis.cls}, would not have been
possible without the following people:
\begin{itemize}
\itemsep 0in
\parsep 0in
\item Francois Pitt --- \programname{ut-thesis.cls} version 1.8, 1999 Dec 10
\item Aichen Low --- original modifications to \programname{ut-thesis.cls} to 
minimally conform to the Georgia Tech requirements.  Also, her
\programname{gt-thesissty.sty} provided additional useful 
code that was incorporated into \programname{gatech-thesis.cls}.
\item Stanford University --- suthesis.sty was adapted by:
\item Ahmed Gheith --- adapted \programname{suthesis.sty} to create 
\programname{GTthesis.sty}.  \programname{GTthesis.sty} was the 
previous ``official'' Georgia Tech style for dissertations.
\item Joonwon Lee, Eilin Tien Lin, and Wei Lui --- modifications to Ahmed's
original \programname{GTthesis.sty}
\item Kalyan Perumalla --- additional modifications to \programname{GTthesis.sty}
(1999 Nov 18)
\item Cody Watson, David Swanson --- even more modifications to 
\programname{GTthesis.sty}
\end{itemize}
\end{acknowledgements}
\end{acknowledgements}
%
\contents
%
\begin{summary}
A brief summary of the entire dissertation should go here.
\end{summary}
\end{preliminary}
%
\chapter{A Shifting Reef}

The year 1866 was signalised by a remarkable incident, a mysterious
and puzzling phenomenon, which doubtless no one has yet forgotten.
Not to mention rumours which agitated the maritime population
and excited the public mind, even in the interior of continents,
seafaring men were particularly excited.  Merchants, common sailors,
captains of vessels, skippers, both of Europe and America,
naval officers\cite{article-full} of all countries, and the 
Governments of several States
on the two continents, were deeply interested in the matter.

For some time past vessels had been met by ``an enormous thing,''
a long object, spindle-shaped, occasionally phosphorescent,
and infinitely larger and more rapid in its movements than a whale.

\section{Add a section marker here}

The facts relating to this apparition (entered in various log-books)
agreed in most respects as to the shape of the object or creature in question,
the untiring rapidity of its movements, its surprising power of locomotion,
and the peculiar life with which it seemed endowed.  If it was a whale,
it surpassed in size all those hitherto classified in science.
Taking into consideration the mean of observations made at divers 
times---rejecting the timid estimate of those who assigned to this object
a length of two hundred feet, equally with the exaggerated opinions
which set it down as a mile in width and three in length---we might fairly
conclude that this mysterious being surpassed greatly all dimensions
admitted by the learned ones of the day, if it existed at all.
And that it DID exist was an undeniable fact; and, with that tendency
which disposes the human mind in favour of the marvellous, we can understand
the excitement produced in the entire world by this supernatural apparition.
As to classing it in the list of fables, the idea was out of the question.

\section{Add another section marker}

\begin{equation}
E=mc^2
\end{equation}

On the 20th of July, 1866, the steamer Governor Higginson,
of the Calcutta and Burnach Steam Navigation Company, had met
this moving mass five miles off the east coast of Australia.
Captain Baker thought at first that he was in the presence of an
unknown sandbank; he even prepared to determine its exact position
when two columns of water, projected by the mysterious object,
shot with a hissing noise a hundred and fifty feet up into the air.
Now, unless the sandbank had been submitted to the intermittent
eruption of a geyser, the Governor Higginson had to do neither
more nor less than with an aquatic mammal, unknown till then,
which threw up from its blow-holes columns of water mixed with
air and vapour.\cite{inbook-full}\cite{book-full}

Similar facts were observed on the 23rd of July in the same year,
in the Pacific Ocean, by the Columbus, of the West India
and Pacific Steam Navigation Company.  But this extraordinary
creature could transport itself from one place to another
with surprising velocity; as, in an interval of three days,
the Governor Higginson and the Columbus had observed it at
two different points of the chart, separated by a distance
of more than seven hundred nautical leagues.

\section{extra section marker}

Fifteen days later, two thousand miles farther off, the Helvetia,
of the Compagnie-Nationale, and the Shannon, of the Royal
Mail Steamship Company, sailing to windward in that portion
of the Atlantic lying between the United States and Europe,
respectively signalled the monster to each other in 42@ 15' N. lat.
and 60@ 35' W. long.  In these simultaneous observations they
thought themselves justified in estimating the minimum length
of the mammal at more than three hundred and fifty feet,
as the Shannon and Helvetia were of smaller dimensions than it,
though they measured three hundred feet over all.

Now the largest whales, those which frequent those parts of the sea round
the Aleutian, Kulammak, and Umgullich islands, have never exceeded the length
of sixty yards, if they attain that.

\begin{equation}
E=mc^2
\end{equation}

In every place of great resort the monster was the fashion.
They sang of it in the cafes, ridiculed it in the papers, and represented
it on the stage.  All kinds of stories were circulated regarding it.
There appeared in the papers caricatures of every gigantic and
imaginary creature, from the white whale, the terrible ``Moby Dick''
of sub-arctic regions, to the immense kraken, whose tentacles could entangle
a ship of five hundred tons and hurry it into the abyss of the ocean.
The legends of ancient times were even revived.

Then burst forth the unending argument between the believers and the
unbelievers in the societies of the wise and the scientific journals.
``The question of the monster'' inflamed all minds.  Editors of
scientific journals, quarrelling with believers in the supernatural,
spilled seas of ink during this memorable campaign, some even drawing blood;
for from the sea-serpent they came to direct personalities.

\section{another section marker}

During the first months of the year 1867 the question seemed buried,
never to revive, when new facts were brought before the public.
It was then no longer a scientific problem to be solved, but a real
danger seriously to be avoided.  The question took quite another shape.
The monster became a small island, a rock, a reef, but a reef of indefinite
and shifting proportions.

\subsection{a subsection marker}

On the 5th of March, 1867, the Moravian, of the Montreal Ocean Company,
finding herself during the night in 27@ 30' lat.  and 72@ 15' long., struck
on her starboard quarter a rock, marked in no chart for that part of the sea.
Under the combined efforts of the wind and its four hundred horse power,
it was going at the rate of thirteen knots.  Had it not been for the superior
strength of the hull of the Moravian, she would have been broken by the shock
and gone down with the 237 passengers she was bringing home from Canada.

The accident happened about five o'clock in the morning, as the day
was breaking.  The officers of the quarter-deck hurried to the after-part
of the vessel.  They examined the sea with the most careful attention.
They saw nothing but a strong eddy about three cables' length distant,
as if the surface had been violently agitated.  The bearings of the place were
taken exactly, and the Moravian continued its route without apparent damage.
Had it struck on a submerged rock, or on an enormous wreck?  They could
not tell; but, on examination of the ship's bottom when undergoing repairs,
it was found that part of her keel was broken.

\subsection{another subsection}

This fact, so grave in itself, might perhaps have been forgotten
like many others if, three weeks after, it had not been re-enacted
under similar circumstances.  But, thanks to the nationality of
the victim of the shock, thanks to the reputation of the company to
which the vessel belonged, the circumstance became extensively circulated.

The 13th of April, 1867, the sea being beautiful, the breeze favourable,
the Scotia, of the Cunard Company's line, found herself in 15@ 12' long.
and 45@ 37' lat.  She was going at the speed of thirteen knots and a half.

At seventeen minutes past four in the afternoon, whilst the passengers were
assembled at lunch in the great saloon, a slight shock was felt on the hull
of the Scotia, on her quarter, a little aft of the port-paddle.

\begin{figure}%
  \centering
  \includegraphics[width=100bp]{jules-verne-test}%
  \caption[A test figure. This is a test figure inserted into this text.  This is a test figure inserted into this text.  This is a test figure inserted into this text.]
  {This is a test figure inserted into this text.}%
  \label{fig:test1}%
\end{figure}

The Scotia had not struck, but she had been struck, and seemingly
by something rather sharp and penetrating than blunt.
The shock had been so slight that no one had been alarmed,
had it not been for the shouts of the carpenter's watch,
who rushed on to the bridge, exclaiming, ``We are sinking! we
are sinking!''  At first the passengers were much frightened,
but Captain Anderson hastened to reassure them.  The danger could
not be imminent.  The Scotia, divided into seven compartments
by strong partitions, could brave with impunity any leak.
Captain Anderson went down immediately into the hold.
He found that the sea was pouring into the fifth compartment;
and the rapidity of the influx proved that the force of the water
was considerable.  Fortunately this compartment did not hold
the boilers, or the fires would have been immediately extinguished.
Captain Anderson ordered the engines to be stopped at once,
and one of the men went down to ascertain the extent of the injury.
Some minutes afterwards they discovered the existence of a
large hole, two yards in diameter, in the ship's bottom.
Such a leak could not be stopped; and the Scotia, her paddles
half submerged, was obliged to continue her course.  She was then
three hundred miles from Cape Clear, and, after three days' delay,
which caused great uneasiness in Liverpool, she entered the basin
of the company.

\subsubsection{a subsubsection with a really really really really really really really really really really long title. I mean a really really really really really really really really really really long title. No, even longer than that.}

The engineers visited the Scotia, which was put in dry dock.
They could scarcely believe it possible; at two yards and a half below
water-mark was a regular rent, in the form of an isosceles triangle.
The broken place in the iron plates was so perfectly defined
that it could not have been more neatly done by a punch.
It was clear, then, that the instrument producing the perforation
was not of a common stamp and, after having been driven with
prodigious strength, and piercing an iron plate 1 3/8 inches thick,
had withdrawn itself by a backward motion.

\subsubsection{a subsubsection with shorter title}

Such was the last fact, which resulted in exciting once more the torrent
of public opinion.  From this moment all unlucky casualties which could
not be otherwise accounted for were put down to the monster.

\section{a final section}

Upon this imaginary creature rested the responsibility of all
these shipwrecks, which unfortunately were considerable;
for of three thousand ships whose loss was annually recorded
at Lloyd's, the number of sailing and steam-ships supposed
to be totally lost, from the absence of all news, amounted to
not less than two hundred!

Now, it was the ``monster'' who, justly or unjustly, was accused
of their disappearance, and, thanks to it, communication between
the different continents became more and more dangerous.
The public demanded sharply that the seas should at any price be
relieved from this formidable cetacean.\footnote{Member of the whale family.}


\chapter{Pro and Con}

At the period when these events took place, I had just returned
from a scientific research in the disagreeable territory
of Nebraska, in the United States.  In virtue of my office
as Assistant Professor in the Museum of Natural History in Paris,
the French Government had attached me to that expedition.
After six months in Nebraska, I arrived in New York towards
the end of March, laden with a precious collection.
My departure for France was fixed for the first days in May.
Meanwhile I was occupying myself in classifying my mineralogical,
botanical, and zoological riches, when the accident happened
to the Scotia.

I was perfectly up in the subject which was the question of the day.
How could I be otherwise?  I had read and reread all the American
and European papers without being any nearer a conclusion.
This mystery puzzled me.  Under the impossibility of forming
an opinion, I jumped from one extreme to the other.
That there really was something could not be doubted,
and the incredulous were invited to put their finger on the wound
of the Scotia.\cite{incollection-full}

\section{An added section}

On my arrival at New York the question was at its height.
The theory of the floating island, and the unapproachable sandbank,
supported by minds little competent to form a judgment, was abandoned.
And, indeed, unless this shoal had a machine in its stomach,
how could it change its position with such astonishing rapidity?

From the same cause, the idea of a floating hull of an enormous
wreck was given up.

\vskip 0.25in
\begin{table}%
 \caption[Add a table]{Add an extra table here. Chinese Menu.}
 \begin{center}
  \begin{tabular}{lp{4.4cm}p{4.4cm}}
                & \textbf{column A} & \textbf{column B} \\
  \textbf{hot}  & Kung Pao Chicken  & General Tso's Chicken \\
  \textbf{mild} & Moo Goo Gai Pan   & Sweet and Sour Pork
  \end{tabular}
  \label{tab:test1}
 \end{center}
\end{table}

There remained, then, only two possible solutions of the question,
which created two distinct parties:  on one side, those who were
for a monster of colossal strength; on the other, those who were
for a submarine vessel of enormous motive power.

\subsection{A subsection}

But this last theory, plausible as it was, could not stand against
inquiries made in both worlds.  That a private gentleman should have
such a machine at his command was not likely.  Where, when, and how
was it built? and how could its construction have been kept secret?
Certainly a Government might possess such a destructive machine.
And in these disastrous times, when the ingenuity of man has
multiplied the power of weapons of war, it was possible that,
without the knowledge of others, a State might try to work such
a formidable engine.

But the idea of a war machine fell before the declaration of Governments.
As public interest was in question, and transatlantic communications
suffered, their veracity could not be doubted.  But how admit that
the construction of this submarine boat had escaped the public eye?
For a private gentleman to keep the secret under such circumstances would
be very difficult, and for a State whose every act is persistently watched
by powerful rivals, certainly impossible.

\begin{figure}%
  \centering%
    \includegraphics[width=100bp]{jules-verne-test}%
  \caption[Another test figure]%
  {This is another test figure inserted into this text
\protect\cite{misc-full}}%
  \label{fig:test2}%
\end{figure}

Upon my arrival in New York several persons did me
the honour of consulting me on the phenomenon in question.
I had published in France a work in quarto, in two volumes,
entitled Mysteries of the Great Submarine Grounds.  This book,
highly approved of in the learned world, gained for me a special
reputation in this rather obscure branch of Natural History.
My advice was asked.  As long as I could deny the reality
of the fact, I confined myself to a decided negative.
But soon, finding myself driven into a corner, I was
obliged to explain myself point by point.  I discussed
the question in all its forms, politically and scientifically;
and I give here an extract from a carefully-studied article
which I published in the number of the 30th of April.
It ran as follows:

\section{Another added section}

``After examining one by one the different theories, rejecting all
other suggestions, it becomes necessary to admit the existence
of a marine animal of enormous power.

\begin{equation}
E=mc^2
\end{equation}

``The great depths of the ocean are entirely unknown to us.
Soundings cannot reach them.  What passes in those remote depths---what 
beings live, or can live, twelve or fifteen miles beneath
the surface of the waters---what is the organisation of these animals,
we can scarcely conjecture.  However, the solution of the problem
submitted to me may modify the form of the dilemma.  Either we do know
all the varieties of beings which people our planet, or we do not.
If we do NOT know them all---if Nature has still secrets in the deeps
for us, nothing is more conformable to reason than to admit the existence
of fishes, or cetaceans of other kinds, or even of new species,
of an organisation formed to inhabit the strata inaccessible to soundings,
and which an accident of some sort has brought at long intervals
to the upper level of the ocean.

``If, on the contrary, we DO know all living kinds, we must
necessarily seek for the animal in question amongst those marine
beings already classed; and, in that case, I should be disposed
to admit the existence of a gigantic narwhal.

``The common narwhal, or unicorn of the sea, often attains
a length of sixty feet.  Increase its size fivefold or tenfold,
give it strength proportionate to its size, lengthen its
destructive weapons, and you obtain the animal required.
It will have the proportions determined by the officers
of the Shannon, the instrument required by the perforation
of the Scotia, and the power necessary to pierce the hull
of the steamer.\cite{inproceedings-full}

``Indeed, the narwhal is armed with a sort of ivory sword,
a halberd, according to the expression of certain naturalists.
The principal tusk has the hardness of steel.  Some of these tusks
have been found buried in the bodies of whales, which the unicorn
always attacks with success.  Others have been drawn out,
not without trouble, from the bottoms of ships, which they
had pierced through and through, as a gimlet pierces a barrel.
The Museum of the Faculty of Medicine of Paris possesses one
of these defensive weapons, two yards and a quarter in length,
and fifteen inches in diameter at the base.

``Very well! suppose this weapon to be six times stronger and the animal
ten times more powerful; launch it at the rate of twenty miles an hour,
and you obtain a shock capable of producing the catastrophe required.
Until further information, therefore, I shall maintain it to be
a sea-unicorn of colossal dimensions, armed not with a halberd,
but with a real spur, as the armoured frigates, or the `rams' of war,
whose massiveness and motive power it would possess at the same time.
Thus may this puzzling phenomenon be explained, unless there be something over
and above all that one has ever conjectured, seen, perceived, or experienced;
which is just within the bounds of possibility.''

\section{The last added section}

These last words were cowardly on my part; but, up to a certain point,
I wished to shelter my dignity as professor, and not give
too much cause for laughter to the Americans, who laugh well
when they do laugh.  I reserved for myself a way of escape.
In effect, however, I admitted the existence of the ``monster.''
My article was warmly discussed, which procured it a high reputation.
It rallied round it a certain number of partisans.  The solution
it proposed gave, at least, full liberty to the imagination.
The human mind delights in grand conceptions of supernatural beings.
And the sea is precisely their best vehicle, the only medium
through which these giants (against which terrestrial animals,
such as elephants or rhinoceroses, are as nothing) can be produced
or developed.

The industrial and commercial papers treated the question chiefly from this
point of view.  The Shipping and Mercantile Gazette, the Lloyd's List,
the Packet-Boat, and the Maritime and Colonial Review, all papers devoted
to insurance companies which threatened to raise their rates of premium,
were unanimous on this point.  Public opinion had been pronounced.
The United States were the first in the field; and in New York they
made preparations for an expedition destined to pursue this narwhal.
A frigate of great speed, the Abraham Lincoln, was put in commission
as soon as possible.  The arsenals were opened to Commander Farragut,
who hastened the arming of his frigate; but, as it always happens,
the moment it was decided to pursue the monster, the monster did not appear.
For two months no one heard it spoken of.  No ship met with it.
It seemed as if this unicorn knew of the plots weaving around it.
It had been so much talked of, even through the Atlantic cable, that jesters
pretended that this slender fly had stopped a telegram on its passage and was
making the most of it.

So when the frigate had been armed for a long campaign, and provided with
formidable fishing apparatus, no one could tell what course to pursue.
Impatience grew apace, when, on the 2nd of July, they learned that a
steamer of the line of San Francisco, from California to Shanghai,
had seen the animal three weeks before in the North Pacific Ocean.
The excitement caused by this news was extreme.  The ship was revictualled
and well stocked with coal.

Three hours before the Abraham Lincoln left Brooklyn pier,
I received a letter worded as follows:

\begin{longquote}
To M. ARONNAX, Professor in the Museum of Paris, Fifth Avenue Hotel, New York.

SIR,--If you will consent to join the Abraham Lincoln
in this expedition, the Government of the United States
will with pleasure see France represented in the enterprise.
Commander Farragut has a cabin at your disposal.

Very cordially yours, J.B. HOBSON, Secretary of Marine.
\end{longquote}

\chapter{I Form My Resolution}

Three seconds before the arrival of J. B. Hobson's letter I no more thought
of pursuing the unicorn than of attempting the passage of the North Sea.
Three seconds after reading the letter of the honourable Secretary of Marine,
I felt that my true vocation, the sole end of my life, was to chase this
disturbing monster and purge it from the world.

But I had just returned from a fatiguing journey, weary and longing
for repose.  I aspired to nothing more than again seeing my country,
my friends, my little lodging by the Jardin des Plantes,
my dear and precious collections---but nothing could keep me back!
I forgot all---fatigue, friends and collections---and accepted without
hesitation the offer of the American Government.

``Besides,'' thought I, ``all roads lead back to Europe; and the unicorn
may be amiable enough to hurry me towards the coast of France.
This worthy animal may allow itself to be caught in the seas of Europe
(for my particular benefit), and I will not bring back less than half
a yard of his ivory halberd to the Museum of Natural History.''
But in the meanwhile I must seek this narwhal in the North
Pacific Ocean, which, to return to France, was taking the road
to the antipodes.\cite{proceedings-full}

``Conseil,'' I called in an impatient voice.

\begin{figure}%
  \centering%
    \includegraphics[width=100bp]{jules-verne-test}%
  \caption[The last test figure]
  {This is the last test figure inserted into this text}%
  \label{fig:test3}%
\end{figure}

Conseil was my servant, a true, devoted Flemish boy, who had accompanied
me in all my travels.  I liked him, and he returned the liking well.
He was quiet by nature, regular from principle, zealous from habit,
evincing little disturbance at the different surprises of life,
very quick with his hands, and apt at any service required of him;
and, despite his name, never giving advice--even when asked for it.

\section{Section marker}

Conseil had followed me for the last ten years wherever science led.
Never once did he complain of the length or fatigue of a journey,
never make an objection to pack his portmanteau for whatever
country it might be, or however far away, whether China or Congo.
Besides all this, he had good health, which defied all sickness,
and solid muscles, but no nerves; good morals are understood.
This boy was thirty years old, and his age to that of his master
as fifteen to twenty.  May I be excused for saying that I was
forty years old?

\begin{equation}
E=mc^2
\end{equation}

But Conseil had one fault:  he was ceremonious to a degree,
and would never speak to me but in the third person,
which was sometimes provoking.

``Conseil,'' said I again, beginning with feverish hands to make
preparations for my departure.

Certainly I was sure of this devoted boy.  As a rule, I never asked
him if it were convenient for him or not to follow me in my travels;
but this time the expedition in question might be prolonged,
and the enterprise might be hazardous in pursuit of an animal capable
of sinking a frigate as easily as a nutshell.  Here there was matter
for reflection even to the most impassive man in the world.
What would Conseil say?

\section{Another marker}

``Conseil,'' I called a third time.

Conseil appeared.

``Did you call, sir?'' said he, entering.

``Yes, my boy; make preparations for me and yourself too.
We leave in two hours.''

``As you please, sir,'' replied Conseil, quietly.

``Not an instant to lose; lock in my trunk all travelling utensils,
coats, shirts, and stockings---without counting, as many as you can,
and make haste.''

``And your collections, sir?'' observed Conseil.

``They will keep them at the hotel.''

``We are not returning to Paris, then?'' said Conseil.

``Oh! certainly,'' I answered, evasively, ``by making a curve.''

``Will the curve please you, sir?''

``Oh! it will be nothing; not quite so direct a road, that is all.
We take our passage in the Abraham, Lincoln.''

``As you think proper, sir,'' coolly replied Conseil.

``You see, my friend, it has to do with the monster---the
famous narwhal.  We are going to purge it from the seas.
A glorious mission, but a dangerous one!  We cannot tell
where we may go; these animals can be very capricious.
But we will go whether or no; we have got a captain who
is pretty wide-awake.''

\section{Yet Another Marker}

Our luggage was transported to the deck of the frigate immediately.
I hastened on board and asked for Commander Farragut.
One of the sailors conducted me to the poop, where I found myself
in the presence of a good-looking officer, who held out his
hand to me.

``Monsieur Pierre Aronnax?'' said he.

``Himself,'' replied I. ``Commander Farragut?''

``You are welcome, Professor; your cabin is ready for you.''

I bowed, and desired to be conducted to the cabin destined for me.

The Abraham Lincoln had been well chosen and equipped
for her new destination.  She was a frigate of great speed,
fitted with high-pressure engines which admitted a pressure
of seven atmospheres.  Under this the Abraham Lincoln attained
the mean speed of nearly eighteen knots and a third an hour---a
considerable speed, but, nevertheless, insufficient to grapple
with this gigantic cetacean.

\vskip 0.25in
\begin{table}%
 \caption[Add another table]{Add another Chinese Menu.}
 \begin{center}%
  \begin{tabular}{lp{4.4cm}p{4.4cm}}
                & \textbf{column A} & \textbf{column B} \\
  \textbf{hot}  & Kung Pao Chicken  & General Tso's Chicken \\
  \textbf{mild} & Moo Goo Gai Pan   & Sweet and Sour Pork
  \end{tabular}
  \label{tab:test2}
 \end{center}
\end{table}

The interior arrangements of the frigate corresponded to its
nautical qualities.  I was well satisfied with my cabin,
which was in the after part, opening upon the gunroom.

``We shall be well off here,'' said I to Conseil.

``As well, by your honour's leave, as a hermit-crab in the shell
of a whelk,'' said Conseil.

I left Conseil to stow our trunks conveniently away, and remounted
the poop in order to survey the preparations for departure.

At that moment Commander Farragut was ordering the last moorings
to be cast loose which held the Abraham Lincoln to the pier
of Brooklyn.  So in a quarter of an hour, perhaps less,
the frigate would have sailed without me.  I should have missed
this extraordinary, supernatural, and incredible expedition,
the recital of which may well meet with some suspicion.

But Commander Farragut would not lose a day nor an hour
in scouring the seas in which the animal had been sighted.
He sent for the engineer.

``Is the steam full on?'' asked he.

``Yes, sir,'' replied the engineer.

``Go ahead,'' cried Commander Farragut.

%
\appendix
%
\chapter{Ned Land}

Captain Farragut was a good seaman, worthy of the frigate he commanded.
His vessel and he were one.  He was the soul of it.  On the question
of the monster there was no doubt in his mind, and he would not allow
the existence of the animal to be disputed on board.  He believed in it,
as certain good women believe in the leviathan--by faith, not by reason.
The monster did exist, and he had sworn to rid the seas of it.  Either Captain
Farragut would kill the narwhal, or the narwhal would kill the captain.
There was no third course.

The officers on board shared the opinion of their chief.
They were ever chatting, discussing, and calculating the various
chances of a meeting, watching narrowly the vast surface of the ocean.
More than one took up his quarters voluntarily in the cross-trees,
who would have cursed such a berth under any other circumstances.
As long as the sun described its daily course, the rigging was
crowded with sailors, whose feet were burnt to such an extent by
the heat of the deck as to render it unbearable; still the Abraham
Lincoln had not yet breasted the suspected waters of the Pacific.
As to the ship's company, they desired nothing better than to meet
the unicorn, to harpoon it, hoist it on board, and despatch it.
They watched the sea with eager attention.

Besides, Captain Farragut had spoken of a certain sum of two thousand dollars,
set apart for whoever should first sight the monster, were he cabin-boy,
common seaman, or officer.

I leave you to judge how eyes were used on board the Abraham Lincoln.

\section{added section marker}

For my own part I was not behind the others, and, left to no one my share
of daily observations.  The frigate might have been called the Argus,
for a hundred reasons.  Only one amongst us, Conseil, seemed to protest
by his indifference against the question which so interested us all,
and seemed to be out of keeping with the general enthusiasm on board.

\begin{equation}
E=mc^2
\end{equation}

I have said that Captain Farragut had carefully provided his
ship with every apparatus for catching the gigantic cetacean.
No whaler had ever been better armed.  We possessed every
known engine, from the harpoon thrown by the hand to the barbed
arrows of the blunderbuss, and the explosive balls of the duck-gun.
On the forecastle lay the perfection of a breech-loading gun,
very thick at the breech, and very narrow in the bore,
the model of which had been in the Exhibition of 1867.
This precious weapon of American origin could throw with ease
a conical projectile of nine pounds to a mean distance
of ten miles.

Thus the Abraham Lincoln wanted for no means of destruction; and, what was
better still she had on board Ned Land, the prince of harpooners.

Ned Land was a Canadian, with an uncommon quickness of hand, and who knew
no equal in his dangerous occupation.  Skill, coolness, audacity, and cunning
he possessed in a superior degree, and it must be a cunning whale to escape
the stroke of his harpoon.

Ned Land was about forty years of age; he was a tall man
(more than six feet high), strongly built, grave and taciturn,
occasionally violent, and very passionate when contradicted.
His person attracted attention, but above all the boldness
of his look, which gave a singular expression to his face.

Who calls himself Canadian calls himself French; and, little communicative
as Ned Land was, I must admit that he took a certain liking for me.
My nationality drew him to me, no doubt.  It was an opportunity for him
to talk, and for me to hear, that old language of Rabelais, which is still
in use in some Canadian provinces.  The harpooner's family was originally
from Quebec, and was already a tribe of hardy fishermen when this town
belonged to France.

Little by little, Ned Land acquired a taste for chatting, and I
loved to hear the recital of his adventures in the polar seas.
He related his fishing, and his combats, with natural poetry
of expression; his recital took the form of an epic poem,
and I seemed to be listening to a Canadian Homer singing the Iliad
of the regions of the North.

\begin{equation}
E=mc^2
\end{equation}

I am portraying this hardy companion as I really knew him.
We are old friends now, united in that unchangeable friendship
which is born and cemented amidst extreme dangers.  Ah, brave Ned!
I ask no more than to live a hundred years longer, that I may have more
time to dwell the longer on your memory.

Now, what was Ned Land's opinion upon the question of the marine monster?
I must admit that he did not believe in the unicorn, and was
the only one on board who did not share that universal conviction.
He even avoided the subject, which I one day thought it my duty
to press upon him.  One magnificent evening, the 30th July (that is
to say, three weeks after our departure), the frigate was abreast
of Cape Blanc, thirty miles to leeward of the coast of Patagonia.
We had crossed the tropic of Capricorn, and the Straits of Magellan
opened less than seven hundred miles to the south.  Before eight
days were over the Abraham Lincoln would be ploughing the waters
of the Pacific.

Seated on the poop, Ned Land and I were chatting of one thing
and another as we looked at this mysterious sea, whose great
depths had up to this time been inaccessible to the eye of man.
I naturally led up the conversation to the giant unicorn, and examined
the various chances of success or failure of the expedition.
But, seeing that Ned Land let me speak without saying too much himself,
I pressed him more closely.

``Well, Ned,'' said I, ``is it possible that you are not convinced
of the existence of this cetacean that we are following?
Have you any particular reason for being so incredulous?''

The harpooner looked at me fixedly for some moments
before answering, struck his broad forehead with his hand
(a habit of his), as if to collect himself, and said at last,
``Perhaps I have, Mr. Aronnax.''

``But, Ned, you, a whaler by profession, familiarised with all
the great marine mammalia---YOU ought to be the last to doubt
under such circumstances!''

``That is just what deceives you, Professor,'' replied Ned.
``As a whaler I have followed many a cetacean, harpooned a great number,
and killed several; but, however strong or well-armed they may
have been, neither their tails nor their weapons would have been
able even to scratch the iron plates of a steamer.''

``But, Ned, they tell of ships which the teeth of the narwhal
have pierced through and through.''

``Wooden ships---that is possible,'' replied the Canadian,
``but I have never seen it done; and, until further proof,
I deny that whales, cetaceans, or sea-unicorns could ever produce
the effect you describe.''

``Well, Ned, I repeat it with a conviction resting on the logic of facts.
I believe in the existence of a mammal power fully organised, belonging to
the branch of vertebrata, like the whales, the cachalots, or the dolphins,
and furnished with a horn of defence of great penetrating power.''

``Hum!'' said the harpooner, shaking his head with the air of a man
who would not be convinced.

``Notice one thing, my worthy Canadian,'' I resumed.
``If such an animal is in existence, if it inhabits the depths
of the ocean, if it frequents the strata lying miles below
the surface of the water, it must necessarily possess an
organisation the strength of which would defy all comparison.''

``And why this powerful organisation?'' demanded Ned.

``Because it requires incalculable strength to keep one's self
in these strata and resist their pressure.  Listen to me.
Let us admit that the pressure of the atmosphere is represented
by the weight of a column of water thirty-two feet high.
In reality the column of water would be shorter, as we are
speaking of sea water, the density of which is greater than
that of fresh water.  Very well, when you dive, Ned, as many
times 32 feet of water as there are above you, so many times
does your body bear a pressure equal to that of the atmosphere,
that is to say, 15 lb.  for each square inch of its surface.
It follows, then, that at 320 feet this pressure equals
that of 10 atmospheres, of 100 atmospheres at 3,200 feet,
and of 1,000 atmospheres at 32,000 feet, that is, about 6 miles;
which is equivalent to saying that if you could attain this
depth in the ocean, each square three-eighths of an inch
of the surface of your body would bear a pressure of 5,600 lb.
Ah! my brave Ned, do you know how many square inches you carry on
the surface of your body?''

``I have no idea, Mr. Aronnax.''

``About 6,500; and as in reality the atmospheric pressure is about 15 lb.
to the square inch, your 6,500 square inches bear at this moment a pressure
of 97,500 lb.''

``Without my perceiving it?''

''Without your perceiving it.  And if you are not crushed by
such a pressure, it is because the air penetrates the interior
of your body with equal pressure.  Hence perfect equilibrium
between the interior and exterior pressure, which thus neutralise
each other, and which allows you to bear it without inconvenience.
But in the water it is another thing.''

``Yes, I understand,'' replied Ned, becoming more attentive;
``because the water surrounds me, but does not penetrate.''

``Precisely, Ned:  so that at 32 feet beneath the surface of the sea you would
undergo a pressure of 97,500 lb.; at 320 feet, ten times that pressure;
at 3,200 feet, a hundred times that pressure; lastly, at 32,000 feet,
a thousand times that pressure would be 97,500,000 lb.---that is to say,
that you would be flattened as if you had been drawn from the plates of
a hydraulic machine!''

``The devil!'' exclaimed Ned.

``Very well, my worthy harpooner, if some vertebrate, several hundred
yards long, and large in proportion, can maintain itself in such depths---of
those whose surface is represented by millions of square inches, that is
by tens of millions of pounds, we must estimate the pressure they undergo.
Consider, then, what must be the resistance of their bony structure,
and the strength of their organisation to withstand such pressure!''

``Why!'' exclaimed Ned Land, ``they must be made of iron plates
eight inches thick, like the armoured frigates.''

``As you say, Ned.  And think what destruction such a mass would cause,
if hurled with the speed of an express train against the hull of a vessel.''

``Yes---certainly---perhaps,'' replied the Canadian, shaken by these figures,
but not yet willing to give in.

``Well, have I convinced you?''

``You have convinced me of one thing, sir, which is that,
if such animals do exist at the bottom of the seas, they must
necessarily be as strong as you say.''

``But if they do not exist, mine obstinate harpooner, how explain
the accident to the Scotia?''

\chapter{At a Venture}

The voyage of the Abraham Lincoln was for a long time marked
by no special incident.  But one circumstance happened which showed
the wonderful dexterity of Ned Land, and proved what confidence
we might place in him.

The 30th of June, the frigate spoke some American whalers,
from whom we learned that they knew nothing about the narwhal.
But one of them, the captain of the Monroe, knowing that Ned Land had
shipped on board the Abraham Lincoln, begged for his help in chasing
a whale they had in sight.  Commander Farragut, desirous of seeing
Ned Land at work, gave him permission to go on board the Monroe.
And fate served our Canadian so well that, instead of one whale,
he harpooned two with a double blow, striking one straight to the heart,
and catching the other after some minutes' pursuit.

Decidedly, if the monster ever had to do with Ned Land's harpoon,
I would not bet in its favour.

The frigate skirted the south-east coast of America with great rapidity.
The 3rd of July we were at the opening of the Straits of Magellan, level with
Cape Vierges.  But Commander Farragut would not take a tortuous passage,
but doubled Cape Horn.

The ship's crew agreed with him.  And certainly it was possible
that they might meet the narwhal in this narrow pass.
Many of the sailors affirmed that the monster could not pass there,
``that he was too big for that!''

The 6th of July, about three o'clock in the afternoon, the Abraham Lincoln,
at fifteen miles to the south, doubled the solitary island,
this lost rock at the extremity of the American continent, to which
some Dutch sailors gave the name of their native town, Cape Horn.
The course was taken towards the north-west, and the next day the screw
of the frigate was at last beating the waters of the Pacific.

``Keep your eyes open!'' called out the sailors.

And they were opened widely.  Both eyes and glasses, a little dazzled,
it is true, by the prospect of two thousand dollars, had not
an instant's repose.

I myself, for whom money had no charms, was not the least
attentive on board.  Giving but few minutes to my meals,
but a few hours to sleep, indifferent to either rain or sunshine,
I did not leave the poop of the vessel.  Now leaning on the netting
of the forecastle, now on the taffrail, I devoured with eagerness
the soft foam which whitened the sea as far as the eye could reach;
and how often have I shared the emotion of the majority of the crew,
when some capricious whale raised its black back above the waves!
The poop of the vessel was crowded on a moment.  The cabins
poured forth a torrent of sailors and officers, each with heaving
breast and troubled eye watching the course of the cetacean.
I looked and looked till I was nearly blind, whilst Conseil kept
repeating in a calm voice:

``If, sir, you would not squint so much, you would see better!''

But vain excitement!  The Abraham Lincoln checked its speed and made
for the animal signalled, a simple whale, or common cachalot,
which soon disappeared amidst a storm of abuse.

But the weather was good.  The voyage was being accomplished under
the most favourable auspices.  It was then the bad season in Australia,
the July of that zone corresponding to our January in Europe,
but the sea was beautiful and easily scanned round a vast circumference.

The 20th of July, the tropic of Capricorn was cut by 105d of longitude,
and the 27th of the same month we crossed the Equator on the 110th meridian.
This passed, the frigate took a more decided westerly direction,
and scoured the central waters of the Pacific.  Commander Farragut thought,
and with reason, that it was better to remain in deep water, and keep
clear of continents or islands, which the beast itself seemed to shun
(perhaps because there was not enough water for him! suggested
the greater part of the crew). The frigate passed at some distance from
the Marquesas and the Sandwich Islands, crossed the tropic of Cancer,
and made for the China Seas.  We were on the theatre of the last diversions
of the monster:  and, to say truth, we no longer LIVED on board.
The entire ship's crew were undergoing a nervous excitement, of which I
can give no idea:  they could not eat, they could not sleep--twenty times
a day, a misconception or an optical illusion of some sailor seated
on the taffrail, would cause dreadful perspirations, and these emotions,
twenty times repeated, kept us in a state of excitement so violent that a
reaction was unavoidable.

And truly, reaction soon showed itself.  For three months,
during which a day seemed an age, the Abraham Lincoln furrowed
all the waters of the Northern Pacific, running at whales,
making sharp deviations from her course, veering suddenly
from one tack to another, stopping suddenly, putting on steam,
and backing ever and anon at the risk of deranging her machinery,
and not one point of the Japanese or American coast
was left unexplored.

The warmest partisans of the enterprise now became its most
ardent detractors.  Reaction mounted from the crew to the captain himself,
and certainly, had it not been for the resolute determination on the part
of Captain Farragut, the frigate would have headed due southward.
This useless search could not last much longer.  The Abraham Lincoln
had nothing to reproach herself with, she had done her best to succeed.
Never had an American ship's crew shown more zeal or patience;
its failure could not be placed to their charge--there remained nothing
but to return.

This was represented to the commander.  The sailors could
not hide their discontent, and the service suffered.
I will not say there was a mutiny on board, but after a reasonable
period of obstinacy, Captain Farragut (as Columbus did)
asked for three days' patience.  If in three days the monster did
not appear, the man at the helm should give three turns of the wheel,
and the Abraham Lincoln would make for the European seas.

This promise was made on the 2nd of November.  It had the effect of
rallying the ship's crew.  The ocean was watched with renewed attention.
Each one wished for a last glance in which to sum up his remembrance.
Glasses were used with feverish activity.  It was a grand defiance
given to the giant narwhal, and he could scarcely fail to answer
the summons and ``appear.''

\begin{equation}
E=mc^2
\end{equation}

Two days passed, the steam was at half pressure; a thousand
schemes were tried to attract the attention and stimulate
the apathy of the animal in case it should be met in those parts.
Large quantities of bacon were trailed in the wake of the ship,
to the great satisfaction (I must say) of the sharks.
Small craft radiated in all directions round the Abraham Lincoln
as she lay to, and did not leave a spot of the sea unexplored.
But the night of the 4th of November arrived without the unveiling of
this submarine mystery.

The next day, the 5th of November, at twelve, the delay would
(morally speaking) expire; after that time, Commander Farragut,
faithful to his promise, was to turn the course to the south-east
and abandon for ever the northern regions of the Pacific.

The frigate was then in 31@ 15' N. lat.  and 136@ 42' E. long.
The coast of Japan still remained less than two hundred miles to leeward.
Night was approaching.  They had just struck eight bells;
large clouds veiled the face of the moon, then in its first quarter.
The sea undulated peaceably under the stern of the vessel.

At that moment I was leaning forward on the starboard netting.
Conseil, standing near me, was looking straight before him.
The crew, perched in the ratlines, examined the horizon which
contracted and darkened by degrees.  Officers with their night
glasses scoured the growing darkness:  sometimes the ocean sparkled
under the rays of the moon, which darted between two clouds,
then all trace of light was lost in the darkness.

In looking at Conseil, I could see he was undergoing a little
of the general influence.  At least I thought so.  Perhaps for
the first time his nerves vibrated to a sentiment of curiosity.

``Come, Conseil,'' said I, ``this is the last chance of pocketing
the two thousand dollars.''

``May I be permitted to say, sir,'' replied Conseil, ``that I never reckoned
on getting the prize; and, had the government of the Union offered a hundred
thousand dollars, it would have been none the poorer.''

``You are right, Conseil.  It is a foolish affair after all, and one upon
which we entered too lightly.  What time lost, what useless emotions!
We should have been back in France six months ago.''

``In your little room, sir,'' replied Conseil, ``and in your museum, sir; and I
should have already classed all your fossils, sir.  And the Babiroussa would
have been installed in its cage in the Jardin des Plantes, and have drawn
all the curious people of the capital!''

``As you say, Conseil.  I fancy we shall run a fair chance of being
laughed at for our pains.''

``That's tolerably certain,'' replied Conseil, quietly; ``I think
they will make fun of you, sir.  And, must I say it---?''

``Go on, my good friend.''

``Well, sir, you will only get your deserts.''

``Indeed!''

``When one has the honour of being a savant as you are, sir, one should
not expose one's self to---''

\begin{equation}
E=mc^2
\end{equation}

Conseil had not time to finish his compliment.
In the midst of general silence a voice had just been heard.
It was the voice of Ned Land shouting:

``Look out there!  The very thing we are looking for---on 
our weather beam!''



%
\begin{postliminary}
\references
\begin{vita}
\begin{vita}
Jules Verne was a really good writer from the Nineteenth century.  He wrote a number of
novels which today would be classified as ``science fiction,'' including \textit{From the
Earth to the Moon} and \textit{Around the World in 80 days}, as well as \textit{20,000 
Leagues Under the Sea}.
\end{vita}

\end{vita}
\end{postliminary}
\end{document}
